\documentclass{article}
\usepackage[utf8]{inputenc}
\usepackage[T1]{fontenc}
\usepackage{geometry}
\geometry{
    a4paper,
    left=2.5cm,
    right=2.5cm,
    top=2.5cm,
    bottom=2.5cm
}
\usepackage{graphicx}
\usepackage{float}
\usepackage{hyperref}
\usepackage[english]{babel}
\usepackage{amsmath}
\usepackage{amssymb}
\usepackage{array}
\usepackage{booktabs}
\usepackage{caption}
\usepackage{subcaption}
\usepackage{enumitem}
\usepackage{fancyhdr}
\usepackage{titlesec}
\usepackage{url}
\usepackage{textcomp}
\usepackage{xcolor}
\usepackage{listings}
\usepackage{multirow}
\usepackage{longtable}
\usepackage{tabularx}
\usepackage{times}

% Set up page style
\pagestyle{fancy}
\fancyhf{}
\rfoot{\thepage}
\renewcommand{\headrulewidth}{0pt}

% Title formatting
\titleformat{\chapter}[display]{\normalfont\bfseries}{\chaptertitle\thechapter}{0pt}{\Huge}
\titlespacing*{\chapter}{0pt}{50pt}{40pt}

% Document info
\title{Template Document}
\author{Academic Support Agent}
\date{\today}

\begin{document}
\maketitle

%----------------------------------------------------------------
% Ficha catalogr fica (OBRIGAT RIO para TCC II)
%  Voc  pode gerar a ficha neste link tamb m: 
%   https://ficha.uffs.edu.br/
%----------------------------------------------------------------
\fichacatalografica

%----------------------------------------------------------------
% Folha de aprova o (OBRIGAT RIO para TCC II)
%   Primeiro par metro   a data da defesa
%   Demais par metros s o os membros da banca (pelo menos 2 nomes)
%   N o precisa passar o nome do Orientador
%----------------------------------------------------------------
\folhadeaprovacao{23/07/2023}{Fulano de Tal - UFFS}{Beltrano de Outro - UFFS}{}


%----------------------------------------------------------------
% Depois da folha de aprova o vem a dedicat ria e a ep grafe.
%----------------------------------------------------------------
\dedicatoria{Dedico este trabalho a meus pais.}

\epigrafe{The art of simplicity is a puzzle of complexity.}
         {Douglas Horton}

%----------------------------------------------------------------
% Tamb m d  para fazer as duas na mesma p gina:
%----------------------------------------------------------------
%\dedigrafe{Dedico este trabalho a meus pais.}
%          {The art of simplicity is a puzzle of complexity.}
%          {Douglas Horton}

%----------------------------------------------------------------
% A seguir, a p gina de agradecimentos (OPCIONAL):
%----------------------------------------------------------------
\begin{agradecimentos}
  lorem ipsum, dolor sit amet consetetur sadipscing elitr sed diam
nonumy eirmod tempor. invidunt ut labore et dolore magna aliquyam

  erad sed, diam voluptua at vero, eos et accusam et justo duo
dolores et ea rebum stet clita.

  kasd gubergren, no sea. takimata sanctus est lorem ipsum dolor sit
amet lorem ipsum dolor sit amet. consetetur sadipscing elitr sed

  diam nonumy, eirmod tempor, invidunt ut labore et dolore magna
aliquyam erat sed diam voluptua at.
\end{agradecimentos}

%----------------------------------------------------------------
% Resumo (OBRIGAT RIO)
%  As palavras chave s o informadas no preambulo l  emcima (\palavraschave)
%----------------------------------------------------------------
\begin{resumo}
Seu resumo em portugu s aqui. lorem ipsum dolor sit amet
consetetur sadipscing elitr sed diam nonumy eirmod tempor invidunt
ut labore et dolore magna aliquyam erat sed diam voluptua at vero
eos et accusam et justo duo dolores et ea rebum stet clita.  kasd
gubergren no sea takimata sanctus est lorem ipsum dolor sit amet
lorem ipsum dolor sit amet consetetur sadipscing elitr sed diam
nonumy eirmod tempor invidunt ut labore et dolore magna aliquyam
erat sed diam voluptua at.
\end{resumo}

%----------------------------------------------------------------
% Abstract (OBRIGAT RIO)
% %  As keywords s o informadas no preambulo l  emcima (\palavraschave)
%----------------------------------------------------------------
\begin{abstract}
Your abstract in English here. lorem ipsum dolor sit amet
consetetur sadipscing elitr sed diam nonumy eirmod tempor invidunt
ut labore et dolore magna aliquyam erat sed diam voluptua at vero
eos et accusam et justo duo dolores et ea rebum stet clita kasd
gubergren no sea takimata sanctus est lorem ipsum dolor sit amet
lorem ipsum dolor sit amet consetetur sadipscing elitr sed diam
nonumy eirmod tempor invidunt ut labore et dolore magna aliquyam
erat sed diam voluptua at
\end{abstract}

%----------------------------------------------------------------
% Listas e sum rio, nessa ordem. Somente o sum rio   obrigat rio,
% portanto, comente as outras listas, caso sejam desnecess rias.
%----------------------------------------------------------------


\listofalgorithms    % Lista de algoritmos   (OPCIONAL)
\listofacronyms      % Lista de siglas       (OPCIONAL)
\listofsymbols       % Lista de s mbolos     (OPCIONAL)


%----------------------------------------------------------------
% Aqui come a o desenvolvimento do trabalho. Para uma melhor
% organiza o do documento, separe-o em arquivos,
% um para cada cap tulo. Para isso, utilize o comando \include,
% como mostrado abaixo.
%----------------------------------------------------------------
% Content from capitulos/capitulo_1 would be included here

% Content from capitulos/capitulo_2 would be included here


%----------------------------------------------------------------
% Aqui vai a bibliografia. 
% Existem 2 estilos de cita o definidos ao instanciar o 
% package abnt2tex: use
% 'alf' para cita es do tipo [Abc+] ou [XYZ] (em ordem
% alfab tica na bibliografia);
% 'num' para cita es
% num ricas do tipo [1], [20], etc., em ordem de refer ncia.
%----------------------------------------------------------------



%----------------------------------------------------------------
% Ap s 
% numera o alfab tica.
%----------------------------------------------------------------

\chapter{Meu primeiro ap ndice}
\chapter{My second appendix}

%----------------------------------------------------------------
% Aqui v o os "cap tulos" de anexos. Cada anexo deve
% ser considerado um cap tulo.
%----------------------------------------------------------------
\anexos
\chapter{Meu primeiro anexo}
\chapter{My second attachment}

% E aqui (para a felicidade de todos) termina o documento.
\end{document}